\documentclass[11pt]{article}
\usepackage[a4paper,margin=1in]{geometry}
\usepackage{amsmath,amssymb,amsthm}
\usepackage{graphicx}
\usepackage{xcolor}
\usepackage[hidelinks]{hyperref}
\usepackage{microtype}
\setlength{\emergencystretch}{2em}

\newtheorem{theorem}{Theorem}
\newtheorem{corollary}[theorem]{Corollary}
\newcommand{\Om}{\Omega}
\newcommand{\T}{\mathbb T}
\newcommand{\R}{\mathbb R}
\newcommand{\Z}{\mathbb Z}
\newcommand{\dd}{\,\mathrm d}

\title{Boundary slope is irrelevant for the channel Poincar\'e constant}
\author{Automated mathematical discovery run \texttt{fa\_banach\_001}}
\date{13 August 2026}

\begin{document}
\maketitle

\begin{center}
\fcolorbox{black}{gray!8}{\parbox{0.91\linewidth}{
\textbf{Status.} Candidate full solution to source Remark~B.2, likely
valid, subject to human review.  The mean-zero Poincar\'e constant of the
periodic channel has an explicit bound independent of the boundary slope.
An additional oscillatory example proves square-root lower growth in the
geometry parameter for the distinct Stokes inf--sup problem.
}}
\end{center}

\section{Source question}

The source is J.~Wilkening, \emph{Inf-sup estimates for the Stokes problem
in a periodic channel}, arXiv:0706.4082 \cite{Source}.  Let
\(\T_L=\R/(L\Z)\) and
\[
 \Om_h=\{(x,y):x\in\T_L,\ 0<y<h(x)\},
 \qquad 0<h_0\le h(x)\le h_1.
\]
The source proves, for mean-zero \(p\in H^1(\Om_h)\),
\(L^{-1}\|p\|_2\le C\|\nabla p\|_2\) with a constant containing the
factor \(1+M\), where \(M=\|h'\|_\infty\).  Remark~B.2 on source
p.~18 asks whether this linear dependence on \(M\) is optimal.  The exact
passage is reproduced below.

\begin{center}
\includegraphics[width=\linewidth]{figures/open_problem_crop.png}
\end{center}

The answer is no: no dependence on \(M\) is needed.  The packet also treats
the separate Remark~4.3, which asks whether quadratic dependence on a
\(C^{1,1}\) geometry parameter is optimal for the Stokes inf--sup inverse.

\section{Proof intuition}

Every channel in the class contains the same flat base rectangle
\(R_0=\T_L\times(0,h_0)\).  Anchor \(p\) to its average on \(R_0\),
where the exact rectangular Poincar\'e estimate is available.  An averaged
one-dimensional identity controls the trace at height \(h_0\).  Values above
that height are then reached by integrating only in the vertical direction.
None of these operations differentiates \(h\), so the boundary slope never
appears.

The Stokes inf--sup constant behaves differently.  A rapidly oscillating
upper boundary creates many upper teeth.  Put alternating signs on those
teeth.  The divergence pairing collapses to an alternating trace functional
on the common interface \(y=h_0\); its \(H^{-1/2}\) norm is of order
\(N^{-1/2}\).  This forces the inf--sup inverse to grow like \(\sqrt N\),
and hence like the square root of the source's geometry parameter.

\section{A slope-free Poincar\'e theorem}

Set
\[
 D=h_1-h_0,\qquad
 \lambda_0=\max\left\{\frac{L}{2\pi},\frac{h_0}{\pi}\right\},
\]
and define
\[
 K_0^2=\max\left\{
 \left(1+\frac{4D}{h_0}\right)\lambda_0^2+\frac{4Dh_0}{3},
 D^2\right\}.
\]

\begin{theorem}[Slope-free mean-zero estimate]\label{thm:poincare}
Let \(h\in C^{0,1}(\T_L)\) satisfy \(h_0\le h\le h_1\).  If
\(p\in H^1(\Om_h)\) and \(\int_{\Om_h}p=0\), then
\[
 \|p\|_{L^2(\Om_h)}\le K_0\|\nabla p\|_{L^2(\Om_h)}.
\]
In particular, \(K_0\) depends only on \(L,h_0,h_1\), not on
\(\|h'\|_\infty\).
\end{theorem}

\begin{proof}
Let \(R_0=\T_L\times(0,h_0)\), let \(c\) be the average of \(p\)
on \(R_0\), and put \(g=p-c\).  Fourier expansion on the periodic
rectangle, with Neumann conditions implicit in the variational inequality,
gives
\begin{equation}\label{eq:rect}
 \|g\|_{L^2(R_0)}\le\lambda_0\|\nabla p\|_{L^2(R_0)}.
\end{equation}
For almost every \(x\), one-dimensional absolute continuity and averaging in
the lower fiber give the exact identity
\[
 g(x,h_0)=\frac1{h_0}\int_0^{h_0}g(x,s)\dd s
 +\frac1{h_0}\int_0^{h_0}t\,p_y(x,t)\dd t.
\]
Consequently,
\begin{equation}\label{eq:trace}
 \int_{\T_L}|g(x,h_0)|^2\dd x
 \le \frac2{h_0}\|g\|_{L^2(R_0)}^2
 +\frac{2h_0}{3}\|p_y\|_{L^2(R_0)}^2.
\end{equation}

Let \(U=\Om_h\setminus R_0\).  For \(h_0<y<h(x)\),
\(g(x,y)=g(x,h_0)+\int_{h_0}^y p_y(x,t)\dd t\).  Integrating in
\(y\), interchanging the order in the second term, and using
\(0\le h(x)-h_0\le D\), we obtain
\begin{equation}\label{eq:upper}
 \|g\|_{L^2(U)}^2
 \le 2D\int_{\T_L}|g(x,h_0)|^2\dd x
 +D^2\|p_y\|_{L^2(U)}^2.
\end{equation}
Equations \eqref{eq:rect}--\eqref{eq:upper} yield
\[
 \|p-c\|_{L^2(\Om_h)}^2
 \le K_0^2\|\nabla p\|_{L^2(\Om_h)}^2.
\]
Finally, the whole-channel mean of \(p\) is zero, so
\[
 \|p-c\|_{L^2(\Om_h)}^2
 =\|p\|_{L^2(\Om_h)}^2+|\Om_h|c^2
 \ge \|p\|_{L^2(\Om_h)}^2.
\]
This proves the theorem.  The fiber calculation for general \(H^1\)
functions follows by the standard density and slicing properties on a
Lipschitz subgraph.
\end{proof}

\begin{corollary}[Answer to Remark B.2]
In the normalization of source Theorem~B.1 one may take
\[
 C=K_0/L,
\]
which is independent of \(M=\|h'\|_\infty\).  Hence the linear
\(M\)-dependence in the source estimate is entirely an artifact of its
boundary-flattening proof.
\end{corollary}

\section{Oscillatory lower bound for the Stokes inf--sup inverse}

For a periodic channel \(\Om\), write
\[
 \beta(\Om)=\inf_{\substack{q\in L^2(\Om),\ \int q=0\\q\ne0}}
 \ \sup_{\substack{\mathbf u\in H^1_0(\Om)^2\\\mathbf u\ne0}}
 \frac{\left|\int_\Om q\,\operatorname{div}\mathbf u\right|}
 {\|q\|_2\|\nabla\mathbf u\|_2}.
\]
Here the zero trace is imposed on the top and bottom walls, with periodicity
in \(x\), as in the source.

\begin{theorem}[Alternating-tooth obstruction]\label{thm:infsup}
Fix \(h_0,D>0\).  For every positive even integer \(N\), let
\[
 h_N(x)=h_0+D(1-\cos Nx),\qquad x\in\T_{2\pi},
\]
and let \(\Om_N=\Om_{h_N}\).  Then
\[
 \beta(\Om_N)\le
 \left(\frac{14\zeta(3)}{\pi^2DN}\right)^{1/2},
 \qquad
 \beta(\Om_N)^{-1}\ge
 \pi\left(\frac{DN}{14\zeta(3)}\right)^{1/2}.
\]
If \(M_N\) is the geometry parameter of source Theorem~4.1, then
\(M_N=\kappa N\) for a fixed \(\kappa=\kappa(h_0,D)>0\).  Thus
\(\beta(\Om_N)^{-1}\ge c(h_0,D)\sqrt{M_N}\).
\end{theorem}

\begin{proof}
The region above \(y=h_0\) has \(N\) teeth
\[
 A_j=\left\{\frac{2\pi j}{N}<x<\frac{2\pi(j+1)}N,
 \ h_0<y<h_N(x)\right\},\qquad 0\le j<N.
\]
Let \(q_N=(-1)^j\) on \(A_j\) and \(q_N=0\) on the base rectangle.
All teeth have the same area and \(N\) is even, hence \(\int q_N=0\),
while
\begin{equation}\label{eq:qnorm}
 \|q_N\|_2^2=\sum_j|A_j|=2\pi D.
\end{equation}

For \(\mathbf u\in H^1_0(\Om_N)^2\), integration on each tooth uses
the zero trace on the curved top.  The horizontal flux integrates to zero at
the endpoints, where the tooth height vanishes, and therefore
\begin{equation}\label{eq:pair}
 -\int_{\Om_N}q_N\operatorname{div}\mathbf u
 =\int_0^{2\pi}w_N(x)u_2(x,h_0)\dd x,
\end{equation}
where \(w_N\) alternates between \(+1\) and \(-1\) on consecutive
cells.  Equivalently, \(w_N(x)=\operatorname{sgn}(\sin(Nx/2))\).

With the Fourier convention
\(\widehat w(k)=(2\pi)^{-1}\int_0^{2\pi}w(x)e^{-ikx}\dd x\), the
square-wave expansion gives nonzero modes only at
\(k=(2m+1)N/2\), and
\begin{equation}\label{eq:fourierw}
 \sum_{k\ne0}\frac{|\widehat w_N(k)|^2}{|k|}
 =\frac{16}{\pi^2N}\sum_{m=0}^\infty\frac1{(2m+1)^3}
 =\frac{14\zeta(3)}{\pi^2N}.
\end{equation}
Put \(v(x)=u_2(x,h_0)\).  For every nonzero horizontal mode,
\[
 \int_0^{h_0}\big(|f'|^2+k^2|f|^2\big)\dd y
 \ge |k|\,|f(h_0)|^2\qquad(f(0)=0),
\]
as follows by expanding \(\int_0^{h_0}|f'-|k|f|^2\).  Parseval
therefore gives
\begin{equation}\label{eq:tracefourier}
 2\pi\sum_{k\ne0}|k||\widehat v(k)|^2
 \le\|\nabla u_2\|_{L^2(\T_{2\pi}\times(0,h_0))}^2
 \le\|\nabla\mathbf u\|_{L^2(\Om_N)}^2.
\end{equation}
Cauchy--Schwarz in \eqref{eq:pair}, followed by
\eqref{eq:fourierw}--\eqref{eq:tracefourier}, yields
\[
 \left|\int_{\Om_N}q_N\operatorname{div}\mathbf u\right|
 \le\left(\frac{28\zeta(3)}{\pi N}\right)^{1/2}
 \|\nabla\mathbf u\|_2.
\]
Divide by \eqref{eq:qnorm} and use \(q_N\) as a competitor in the
definition of \(\beta\).

Finally, \(\|h_N'\|_\infty=DN\), and
\(\frac12\|h_Nh_N''\|_\infty\) is a fixed positive multiple of
\(N^2\).  Hence the source parameter
\[
 M_N^2=\max\left\{\|h_N'\|_\infty^2,
 \frac12\|h_Nh_N''\|_\infty\right\}
\]
equals \(\kappa^2N^2\) for fixed \(\kappa>0\).
\end{proof}

\section{Verification, limitations, and novelty}

Theorem~\ref{thm:poincare} fully answers the extracted Remark~B.2.  It
does not claim the displayed constant is sharp in \(L,h_0,h_1\).  Theorem
\ref{thm:infsup} is a substantial partial result for the distinct Remark~4.3:
the source supplies an \(O(M^2)\) upper bound, while the new family forces
\(\Omega(\sqrt M)\) growth.  The exponent between \(1/2\) and \(2\)
remains open.

Eight focused checks and upgrade attempts covered the trace coefficient, the
two different means, Sobolev slicing, exact square-wave normalization,
zero-mode removal, tooth flux identity, a direct vertical divergence
construction, Piola flattening, and quantitative Bogovski\u\i{}/John-domain
estimates.  The latter routes reintroduced uncontrolled horizontal
derivatives or curvature and did not justify a universal subquadratic upper
bound.

A bounded novelty search used the exact title, arXiv id, author, both quoted
optimality phrases, and combinations of ``periodic channel'', ``Poincar\'e
constant'', ``LBB constant'', ``slope'', and ``oscillating boundary''.  It
also checked general uniform-Poincar\'e work \cite{Ruiz} and quantitative
Bogovski\u\i{} estimates \cite{Duran,GuzmanSalgado}.  No explicit answer to
Remark~B.2, no base-rectangle proof, and no alternating-tooth lower bound for
Remark~4.3 were found.  Novelty confidence is provisional pending expert
review.

\begin{thebibliography}{9}
\bibitem{Source}
J.~Wilkening,
\emph{Inf-sup estimates for the Stokes problem in a periodic channel},
arXiv:0706.4082 (2007), especially Remarks~4.3 and B.2.

\bibitem{Duran}
R.~G.~Dur\'an,
\emph{An elementary proof of the continuity from $L^2_0(\Omega)$ to
$H^1_0(\Omega)^n$ of Bogovski\u\i{}'s right inverse of the divergence},
arXiv:1103.3718 (2011).

\bibitem{Ruiz}
D.~Ruiz,
\emph{A note on the uniformity of the constant in the Poincar\'e
inequality}, arXiv:1208.6045 (2012).

\bibitem{GuzmanSalgado}
J.~Guzm\'an and A.~J.~Salgado,
\emph{Estimation of the continuity constants for Bogovski\u\i{} and
regularized Poincar\'e integral operators}, arXiv:2010.04105 (2020).
\end{thebibliography}

\end{document}
